\section{Objetivos de seguridad y protección}

\subsection{Los tres elementos de la CIA}

\begin{importantbox}{Objetivos de seguridad de los sistemas Agentic: CIA}
\begin{itemize}
    \item \textbf{Confidentiality (confidencialidad)}: la información solo es accesible para entidades autorizadas, incluidos los secretos del sistema, las credenciales de usuario, los datos de usuario y el propio modelo
    \item \textbf{Integrity (integridad)}: los sistemas y los datos no han sido alterados, y se mantienen exactos y confiables
    \item \textbf{Availability (disponibilidad)}: los usuarios autorizados acceden de manera confiable y oportuna a los datos, los sistemas y los recursos
\end{itemize}
\end{importantbox}

\subsection{Comparación con los sistemas tradicionales}

En comparación con los sistemas de cómputo tradicionales, los sistemas de Agentic AI tienen \textbf{objetivos de protección adicionales}:
\begin{itemize}
    \item \textbf{Confidencialidad}: se necesita proteger de manera adicional las claves de API, los prompts secretos, el historial de interacciones, los parámetros propietarios del modelo, entre otros
    \item \textbf{Integridad}: se necesita garantizar la integridad del modelo (model integrity)
    \item \textbf{Disponibilidad}: se necesita garantizar el rendimiento del modelo y la disponibilidad del servicio
\end{itemize}

\begin{warningbox}{El aumento de la superficie de ataque introducido por los LLM}
El uso de LLM amplía de manera significativa la superficie de ataque:
\begin{itemize}
    \item Confidencialidad: después de procesar información sensible, el LLM puede filtrarla a través de su salida
    \item Integridad: el prompt puede aprovechar datos no confiables; el modelo puede ser envenenado o contener puertas traseras
    \item Disponibilidad: el propio modelo puede sufrir ataques de denegación de servicio
\end{itemize}
\end{warningbox}

